%\errorcontextlines 10000
\documentclass[aspectratio=169]{beamer}

\usetheme{Berlin}
\setbeamertemplate{navigation symbols}{}
\setbeamertemplate{headline}{}
\setbeamertemplate{footline}{}

\defbeamertemplate{section page}{nonumber}[1][]{%
  \begin{centering}
    {\usebeamerfont{section name}\usebeamercolor[fg]{section name}#1}
    \vskip1em\par
    \begin{beamercolorbox}[sep=12pt,center]{part title}
      \usebeamerfont{section title}\insertsection\par
    \end{beamercolorbox}
  \end{centering}
}

\usepackage{braket}
\usepackage{esdiff}

\newcommand{\sqrtpm}{\sqrt{p^2 + m^2}}
\newcommand{\envelope}{\frac{p}{\sqrtpm}}
\newcommand{\integrand}{\frac{p e^{ipr}}{\sqrt{p^2 + m^2}}}
\newcommand{\CC}{\mathbb{C}}

\newcommand{\vp}{\mathbf{p}}
\newcommand{\vx}{\mathbf{x}}
\newcommand{\vy}{\mathbf{y}}
\newcommand{\vr}{\mathbf{r}}
\newcommand{\vpp}{\mathbf{p}'}
\newcommand{\omp}{\omega_{p}}
\newcommand{\ommp}{\omega_{-p}}
\newcommand{\omvp}{\omega_{\vp}}
\newcommand{\omvpp}{\omega_{\vpp}}
\newcommand{\ap}{a_{\vp}}
\newcommand{\app}{a_{\vpp}}
\newcommand{\adp}{a^\dagger_{\vp}}
\newcommand{\adpp}{a^\dagger_{\vpp}}

\newcommand{\di}[1]{\diff{}{#1}}
\newcommand{\diffop}{\mathop{}\!\mathrm{d}}
\newcommand{\dip}{\diffop p}
\newcommand{\du}{\diffop u}
\newcommand{\dxp}{\diffop x'}
\newcommand{\dyp}{\diffop y'}
\newcommand{\dxi}{\diffop \xi}
\newcommand{\drp}{\diffop r'}
\newcommand{\drho}{\diffop \rho}
\newcommand{\dtheta}{\diffop \theta}

\newcommand{\bysqrt}[1]{\frac{1}{\sqrt{#1}}}
\newcommand{\intpthree}{\int \frac{\diffop^3\vp}{(2\pi)^3}}
\newcommand{\intpthreeppthree}{\int \frac{\diffop^3\vp}{(2\pi)^3} \frac{\diffop^3\vpp}{(2\pi)^3}}
\newcommand{\com}[2]{[{#1},{#2}]}
\newcommand{\phixy}{\Braket{0|\phi(x)\phi(y)|0}}
\newcommand{\intii}{\int_{-\infty}^{+\infty}}
\newcommand{\limdel}{\lim_{\delta \to 0}}
\newcommand{\integral}{\intii \dip \integrand}
\newcommand{\cintegral}{\int_m^\infty \drho \frac{\rho e^{-\rho r}}{\sqrt{\rho^2 - m^2}}}

\newcommand{\phit}{\tilde{\phi}}
\newcommand{\sigmat}{\tilde{\sigma}}

\renewcommand\Re{\operatorname{Re}}
\renewcommand\Im{\operatorname{Im}}

\title{Contour integral for the free field propagator}
\author{Edwin Steiner}
\date{\today}
% XXX add git hash
% XXX credits
% XXX GPL license for code
% XXX blink branch cut?
% XXX link to github code


\begin{document}

%: $opts = '--n 1250 --z0 0.1 --z1 10 --m 1';

%: seq_name('teaser');
%: animate(3, sub {
%:     my $Pr = interpol_xaby(1, 0.1, 0.1, 0.1, 15, 70, 70, 70, $_[0]);
%:     [sprintf($opts.' --Pr %.2g --Pi %.2g --d 0 --smear 0.0333333', $Pr, $Pr),
%:      sprintf(qq{
%:          --pre-plot 'set label "1/σ = 30" at 5,8 font "serif,26"'
%:          --pre-plot 'set label "Λ = %.2g" at 5,5 font "serif,26"'}, $Pr),
%:      $im_ints]
%: });


\begin{frame}
\titlepage
\begin{tiny}
\includegraphics[width=0.1\textwidth]{CC_BY.png}

\medskip
This work is licensed under a Creative Commons Attribution 4.0 International License.

\url{http://creativecommons.org/licenses/by/4.0/}
\end{tiny}
\end{frame}


\section{Introduction}


\begin{frame}
For a free massive \alert{scalar field}
%.2

\pause
$$\phi(x) = \intpthree \frac{1}{\sqrt{2\omvp}} \left\{ e^{-ipx} \ap + e^{ipx} \adp \right\}$$

with

$$\omvp = \sqrt{\vp^2 + m^2}$$

\pause
we want to calculate the \alert{2-point function}
%.2

\pause
$$\phixy$$
%.2

\pause
\medskip
for \alert{purely spatial} separations $\vx - \vy = \vr, x^0 - y^0 = 0$.
%.4
\end{frame}


\begin{frame}
Plugging in the fields and using $\com{\ap}{\adpp} = (2\pi)^3 \delta(\vp - \vpp)$ we get:

\pause
$$\phixy = \frac{-i}{2(2\pi)^2 r} \intii \frac{p e^{ipr}}{\sqrtpm} \dip$$

\pause
\bigskip
(See end of video for details.)
%.1

\pause
\medskip
See also Peskin \& Schroeder, \textit{An Introduction to Quantum Field Theory}, 2.50 to 2.52.
%.4
\end{frame}


\begin{frame}
The only step left is to evaluate the integral
%.2

\pause
$$\int_{-\infty}^{\infty} \dip \integrand$$

\pause
\medskip
\alert{This integral does not converge!}

\pause
\bigskip
\bigskip
Riemann -- ``$\lim_{b \to \infty} \int_0^b \ldots$ does not exist.'' (It oscillates.)

\pause
\bigskip
Lebesgue -- ``Positive and negative $\Re$ / $\Im$ parts are bounded away from zero on sets of infinite measure.''
%.6
\end{frame}


\section{The integrand}

\begin{frame}
Let's consider the integrand

$$\integrand$$
%.2

\pause
as a function of a \alert{complex} variable $p$.
\end{frame}


\begin{frame}
Actually, let's first look at:
%.1

\pause
$$\envelope$$

\pause
No surprises on the real line:
%.2

\pause
$$\envelope \ \alert{\rightarrow \ 1} \ \textrm{ for } p \rightarrow +\infty,$$
%.1

\pause
$$\envelope \ \alert{\rightarrow \ -1} \ \textrm{ for } p \rightarrow -\infty.$$

\end{frame}


%: animate(3, sub { ['--envelope --z0 -10 --z1 10 --m 1', '', $im_env] });

%: animate(10, sub {
%:     my $m = interpol_xaby(5, 0.1, 0.1, 0.1, 1, 0.2, 3, 1, $_[0]);
%:     [sprintf('--envelope --z0 -10 --z1 10 --m %.4g', $m),
%:      sprintf(qq{
%:          --pre-plot 'set label "varying the mass..." at 0.25,2 font "serif,26"'
%:          --pre-plot 'set label "m = %4.2g" at 0.25,1.5 font "serif,26"'}, $m),
%:      $im_env]
%: });


\begin{frame}
It gets more interesting for complex $p$:

\pause
$$\envelope$$
%..5

\begin{itemize}

\pause
\item For $|p| \rightarrow \infty$ this goes to \alert{$\pm 1$} in the right/left half plane.

\pause
\item There are branch point \alert{singularities} at $p = im$ and $p = -im$.

\pause
\item The \alert{branch cut} lies along the imaginary axis towards $+i\infty$ and $-i\infty$.
%.3
\pause

(From branch cut of $\sqrt{\ldots}$ on the negative real axis.)
%.6
\end{itemize}
\end{frame}


%: seq_name('branch-cut');
%: animate(20, sub {
%:     my $y = interpol_xaby(4, 0.1, 0.1, 0.1, 0, 2, 0, 0, $_[0]);
%:     my $preplot = '';
%:     my $im = '';
%:     if ($y >= 1.0) {
%:         $preplot .= q{ --pre-plot 'set grid x2tics ; set x2tics (0.5) format "" scale 0'};
%:     }
%:     if ($y > 1.2) {
%:         $preplot .= q{ --pre-plot 'set label "branch cut" at 0.5,2 center font "serif,26"'};
%:     }
%:     if ($y > 1.4) {
%:         $preplot .= q{ --pre-plot 'set arrow from 0.45,1.5 to 0.49,1.5 head'};
%:         $preplot .= q{ --pre-plot 'set arrow from 0.55,1.5 to 0.51,1.5 head'};
%:         $im .= ' label-branch-cut-neg.png -geometry +358+150 -composite';
%:         $im .= ' label-branch-cut-pos.png -geometry +730+150 -composite';
%:     }
%:     [sprintf('--envelope --z0 -10,%.4g --z1 10,%.4g --m 1', $y, $y),
%:      $preplot, $im_env.$im]
%: });


\begin{frame}
Including the exponential we get
%.1

\pause
$$\integrand$$

which is
%.1

\begin{itemize}
\pause
\item \alert{oscillating} along the real axis,

\pause
\item \alert{exponentially suppressed} in the upper half plane.
%.4
\end{itemize}
\end{frame}


%: animate(3, sub { ['--integrand --z0 -10 --z1 10 --m 1 --r 2', '', $im_igd] });

%: animate(12, sub {
%:     my $r = interpol_xaby(5, 0.1, 0.1, 0.1, 2, 0.5, 10, 2,$_[0]);
%:     [sprintf('--integrand --z0 -10 --z1 10 --m 1 --r %.3g', $r),
%:      sprintf(qq{
%:          --pre-plot 'set label "varying r..." at 0.25,2 font "serif,26"'
%:          --pre-plot 'set label "r = %4.2g" at 0.25,1.5 font "serif,26"'}, $r),
%:      $im_igd]
%: });

%: animate(10, sub {
%:     my $y = interpol_xaby(2, 0.1, 0.1, 0.1, 0, 2, 2, 0, $_[0]);
%:     my $preplot = ($y > 1.2) ? 'set label "branch cut" at 0.5,2 center font "serif,26"' : '';
%:     [sprintf('--integrand --z0 -10,%.3g --z1 10,%.3g --m 1', $y, $y),
%:      "--pre-plot '$preplot'", $im_igd]
%: });

%: seq_name('rotate-1');
%: animate(8, sub {
%:     my $t = interpol_xaby(1, 0.1, 0.2, 0.2, 0, 0, 1, 1, $_[0]);
%:     my @z0 = interpol_complex(-10,0, 0.01,0, $t);
%:     my @z1 = interpol_complex(+10,0, 0.01,2, $t);
%:     [sprintf('--integrand --z0 %.3g,%.3g --z1 %.3g,%.3g --m 1 --r 2', @z0, @z1),
%:      '', $im_igd]
%: });
%: seq_name('rotate-2');
%: animate(8, sub {
%:     my $t = interpol_xaby(2, 0.1, 0.1, 0.1, 0, 1, 1, 0, $_[0]);
%:     my @z0 = interpol_complex(0.01,0, 8,0, $t);
%:     my @z1 = interpol_complex(0.01,2, 8,2, $t);
%:     [sprintf('--integrand --z0 %.3g,%.3g --z1 %.3g,%.3g --m 1 --r 2', @z0, @z1),
%:      '', $im_igd]
%: });
%: seq_name('rotate-3');
%: animate(8, sub {
%:     my $t = interpol_xaby(1, 0.1, 0.2, 0.2, 1, 1, 0, 0, $_[0]);
%:     my @z0 = interpol_complex(-10,0, 0.01,0, $t);
%:     my @z1 = interpol_complex(+10,0, 0.01,2, $t);
%:     [sprintf('--integrand --z0 %.3g,%.3g --z1 %.3g,%.3g --m 1 --r 2', @z0, @z1),
%:      '', $im_igd]
%: });


\section{The integral}

\begin{frame}
Let's do some numerical integration!
\end{frame}


\begin{frame}
We try a sharp \alert{cutoff} at $p = \pm \Lambda$:
%.2

\pause
$$\int_{-\Lambda}^{+\Lambda} \dip \integrand$$

\pause
\bigskip
It looks like this \alert{as a function of $r$}...
\end{frame}


%: animate(3, sub {
%:     [$opts.' --Pr 45 --Pi 0 --d 0', '', $im_int]
%: });

%: animate(20, sub {
%:     my $Pr = interpol_xaby(5, 0.1, 0.1, 0.1, 45, 25, 60, 45, $_[0]);
%:     [sprintf($opts.' --Pr %.3g --Pi 0 --d 0', $Pr),
%:      sprintf(qq{
%:          --pre-plot 'set label "varying the cutoff..." at 5,8 font "serif,26"'
%:          --pre-plot 'set label "Λ = %.2g" at 5,5 font "serif,26"'}, $Pr),
%:      $im_int]
%: });


\begin{frame}
We see:
%.1

\begin{itemize}
\pause
\item strong \alert{oscillations} at the cutoff momentum,

\pause
\item \alert{``average value''} falling \alert{exponentially} with $r$.
\end{itemize}

\pause
\bigskip
By deforming the contour we can show:

\pause
\begin{itemize}
\item \alert{only} the cutoff $p = \Lambda$ creates the oscillations.
\end{itemize}
%.6
\end{frame}


\begin{frame}
Witness the miracle of \alert{Cauchy's theorem}!
\end{frame}


%: animate(3, sub { [$opts.' --Pr 45 --Pi 0 --d 0', '', $im_int] });

%: seq_name('cauchy');
%: animate(10, sub {
%:     my $t = interpol_xaby(1, 0.1, 0.1, 0.1, 0, 45, 0, 45, $_[0]);
%:     my $preplot = q{--pre-plot 'set label "pushing the contour..." at 5,8 font "serif,26"'};
%:     if ($_[0] > 0.1) {
%:         $preplot .= q{ --pre-plot 'set arrow from 1.9,6.8 to 1,5 head'};
%:         $preplot .= q{ --pre-plot 'set label "no change" at 2,7 font "serif,18"'};
%:     }
%:     [sprintf($opts.' --Pr 45 --Pi %.3g --d 0', $t),
%:      $preplot, $im_int]
%: });

%: animate(10, sub {
%:     my $Pr = interpol_xaby(3, 0.1, 0.1, 0.1, 45, 25, 60, 45, $_[0]);
%:     [sprintf($opts.' --Pr %.3g --Pi 45 --d 0', $Pr),
%:      sprintf(qq{
%:          --pre-plot 'set label "varying the cutoff..." at 5,8 font "serif,26"'
%:          --pre-plot 'set label "Λ = %.2g" at 5,5 font "serif,26"'}, $Pr),
%:      $im_int]
%: });


\begin{frame}
Parts of the contour...
%.1

\pause
\begin{itemize}
\item \alert<2>{... at large imaginary values are negligible}

%: seq_name('contour-top');
%: animate(10, sub {
%:     my $t = $_[0];
%:     my $preplot = '';
%:     if ($_[0] > 0.1) {
%:         $preplot .= q{ --pre-plot 'set arrow from 1.9,6.8 to 1,5 head'};
%:         $preplot .= q{ --pre-plot 'set label "no (visible) change" at 2,7 font "serif,18"'};
%:     }
%:     if ($_[0] > 0.2) {
%:         $preplot .= q{ --pre-plot 'set label "omitting top part" at 2.5,-6.5 font "serif,16" textcolor rgb "red"'};
%:     }
%:     my $copts = '';
%:     if ((($t >= 0.25) && ($t < 0.5)) || ($t >= 0.75)) {
%:         $copts = ' --contour-IUI';
%:     }
%:     [sprintf($opts.' --Pr 45 --Pi 45 --d 0'.$copts),
%:      $preplot, $im_int]
%: });

\pause
\item \alert<3>{... at $\Re p = \pm \Lambda$ contribute the cutoff oscillations}

%: seq_name('contour-ii');
%: animate(10, sub {
%:     my $t = $_[0];
%:     my $preplot = '';
%:     my $copts = ' --contour-IUI';
%:     if ($t > 0.1) {
%:         $preplot .= q{ --pre-plot 'set label "omitting middle part" at 2.3,-6.5 font "serif,16" textcolor rgb "red"'};
%:     }
%:     if ((($t >= 0.25) && ($t < 0.5)) || ($t >= 0.75)) {
%:         $copts = ' --contour-II';
%:     }
%:     [sprintf($opts.' --Pr 45 --Pi 45 --d 0'.$copts),
%:      $preplot, $im_int]
%: });

\pause
\item \alert<4>{... close to $\Re p = 0$ contribute the exponential part}

%: seq_name('contour-u');
%: animate(10, sub {
%:     my $d = interpol_xaby(1, 0.1, 0.1, 0.7, 0, 1, 1, 1, $_[0]);
%:     my $preplot = '';
%:     [sprintf($opts.' --Pr 45 --Pi 45 --d %.2g --contour-IUI', $d),
%:      $preplot, $im_int]
%: });

\end{itemize}
\end{frame}


\begin{frame}
Apparently the oscillations are an \alert{artifact of the cutoff}.

\bigskip
\pause
More carefully:
%.1

\pause
\medskip
We consider the fields as \alert{distributions}...
%.2

\pause
...and \alert{smear} them with test functions.
%.3
\end{frame}


\begin{frame}
For simplicity we try:
%.1

\pause
\alert{Gaussian} smearing $e^{-\frac{(x' - x)^2}{2\sigma^2}}$ of $\phi$ only along $\vr$.
%.5

\medskip
\pause
This causes an \alert{exponential suppression of high momenta} $p$:
%.2

\pause
\begin{equation*}
\intii \dip \integrand \alert{e^{-(\sigma p)^2}}
\end{equation*}
%.1

\pause
\medskip
(See end of video for details.)
\end{frame}


%: seq_name('smear-integrand');
%: animate(12, sub {
%:     my $sigma = interpol_xaby(3, 0.2, 0.2, 0.1, 0, 0.2, 0.01, 1/30, $_[0]);
%:     [sprintf('--integrand --z0 -10 --z1 10 --m 1 --r 2 --smear %.2g', $sigma),
%:      sprintf(qq{
%:          --pre-plot 'set label "varying σ..." at 0.25,2 font "serif,26"'
%:          --pre-plot 'set label "σ = %4.2g" at 0.25,1.5 font "serif,26"'}, $sigma),
%:      $im_igds]
%: });


\begin{frame}
Let's smear the fields by $\sigma$...

\medskip
\pause
...and then push the contour to $\Lambda \gg \frac{1}{\sigma}$.
\end{frame}


%: seq_name('smear-integral');
%: animate(24, sub {
%:     my $Pr = interpol_xaby(4, 0.1, 0.1, 0.1, 15, 70, 15, 70, $_[0]);
%:     [sprintf($opts.' --Pr %.2g --Pi %.2g --d 0 --smear 0.0333333', $Pr, $Pr),
%:      sprintf(qq{
%:          --pre-plot 'set label "1/σ = 30" at 5,8 font "serif,26"'
%:          --pre-plot 'set label "Λ = %.2g" at 5,5 font "serif,26"'}, $Pr),
%:      $im_ints]
%: });


\begin{frame}
All the ``outer'' parts of the contour are exponentially small now.

\medskip
\pause
Let's neglect them...
\end{frame}


%: seq_name('smear-integral-2');
%: animate(12, sub {
%:     my $d = interpol_xaby(1, 0.2, 0.1, 0.2, 0, 2, 2, 2, $_[0]);
%:     [sprintf($opts.' --Pr 70 --Pi 70 --d %.2g --smear 0.0333333', $d),
%:      sprintf(qq{
%:          --pre-plot 'set label "1/σ = 30" at 5,8 font "serif,26"'
%:          --pre-plot 'set label "Λ = %.2g" at 5,5 font "serif,26"'}, 70),
%:      $im_ints]
%: });


\begin{frame}
For this modified contour, smearing makes no difference (in the limit)...
\end{frame}


%: seq_name('unsmear-integral-2');
%: animate(6, sub {
%:     my $isigma = interpol_xaby(1, 0.2, 0.1, 0.2, 30, 1000, 1000, 1000, $_[0]);
%:     my $smear;
%:     my $text;
%:     my $im;
%:     if ($isigma <= 999) {
%:         $smear = sprintf(' --smear %.3g', 1/$isigma);
%:         $text = sprintf('1/sigma = %.4g', $isigma);
%:         $im = $im_ints;
%:     }
%:     else {
%:         $smear = '';
%:         $text = 'not smeared';
%:         $im = $im_int;
%:     }
%:     [$opts.' --Pr 70 --Pi 70 --d 2'.$smear,
%:      sprintf(qq{
%:          --pre-plot 'set label "$text" at 5,8 font "serif,26"'
%:          --pre-plot 'set label "Λ = %.2g" at 5,5 font "serif,26"'}, 70),
%:      $im]
%: });


\begin{frame}
In the limit, the integral along this \alert{U}-shaped contour $C$ becomes
%.3

\pause
\begin{equation*}
\int_C \dip \integrand = 2i \alert{\cintegral}
\end{equation*}

\medskip
(See end of video for details.)
%.6
\end{frame}


\begin{frame}
This integral \alert{converges} and we get \alert{approximately}
%.3

\pause
\begin{equation*}
\phixy \approx \sqrt{\frac{m}{32\pi^3 r^3}} \ \alert{e^{-mr}}
\end{equation*}

\medskip
for $r \gg \frac{1}{m}$

\medskip
(See end of video for details.)
%.6
\end{frame}


\begin{frame}
The \alert{exact} result can be written with a \alert{modified Bessel function} $\mathrm{K}_1$:
%.3

\pause
\begin{equation*}
\phixy = \frac{m}{(2\pi)^2 r} \  \alert{\mathrm{K}_1(mr)}
\end{equation*}

\medskip
(See end of video for details.)
%.6
\end{frame}


\begin{frame}
Let's \alert{compare}:
%.1

\begin{enumerate}
\pause
\item our \alert<2>{numerical} result with
%.2
\pause
\item \alert<3>{exact} and
%..5
\pause
\item \alert<4>{approximate} analytic results...
%.4
\end{enumerate}
\end{frame}


%: seq_name('approx');
%: animate(20, sub {
%:     my $t = $_[0];
%:     my $y = interpol_xaby(1, 0.4, 0.1, 0.5, 15, 0.1, 0.1, 0.1, $t);
%:     my $plot = '';
%:     my $text = 'numerical';
%:     if ($t > 0.1) {
%:         $plot .= qq{ --append-plot '"${opt_prefix}bessel-FUNCTION.dat" using 2:5 with lines lc rgb "black" lt 2 lw 3 title "2*m*K1(mr)"' };
%:         $text = 'exact';
%:     }
%:     if ($t > 0.2) {
%:         $plot .= q{ --pre-plot 'm = 1' --append-plot 'sqrt(2*pi*m/x)*exp(-m*x) with lines lc rgb "black" lt 3' };
%:         $text = 'exact & approximate';
%:     }
%:     if ($t > 0.3) {
%:         $text = 'zooming...';
%:     }
%:     $plot .= qq{ --pre-plot 'set label "$text" at 5,8 font "serif,26"' };
%:     [sprintf($opts.' --Pr 70 --Pi 70 --d 2'),
%:      sprintf(' --yrange %.2g:%.2g', -$y, $y).$plot,
%:      $im_int]
%: });


\section{Credits}

\begin{frame}
\frametitle{Software Credits}

Numerical integration: GNU Scientific Library

Plotting: gnuplot

Typesetting: \TeX, \LaTeX, Beamer

Conversion: ImageMagick, avconv
%.6
\end{frame}


\section{Calculational details}

\begin{frame}
\setbeamertemplate{section page}[nonumber]
\sectionpage
\end{frame}

\begin{frame}
Plugging in the fields, we get:

\begin{equation*}
\begin{split}
& \phixy = \intpthreeppthree \bysqrt{2\omvp} \bysqrt{2\omvpp} \\
\bigg\{ & e^{-ipx-ip'y} \braket{0|\ap\app|0} + e^{-ipx+ip'y} \braket{0|\ap\adpp|0} \\
        & + e^{ipx-ip'y} \braket{0|\adp\app|0} + e^{ipx+ip'y} \braket{0|\adp\adpp|0} \bigg\}
\end{split}
\end{equation*}

using $\ap\ket{0} = 0$:

\begin{equation*}
= \intpthreeppthree \bysqrt{2\omvp} \bysqrt{2\omvpp}
\bigg\{ e^{-ipx+ip'y} \braket{0|\ap\adpp - \adpp\ap|0} \bigg\}
\end{equation*}
\end{frame}


\begin{frame}
using $\com{\ap}{\adpp} = (2\pi)^3 \delta(\vp - \vpp)$, integrating out $\vpp$:

\begin{equation*}
\phixy = \intpthree \frac{1}{2\omvp} e^{-ip(x-y)}
\end{equation*}

for $x^0 - y^0 = 0$:

\begin{equation*}
\phixy = \intpthree \frac{1}{2\omvp} e^{i\vp\cdot\vr}
\end{equation*}
\end{frame}


\begin{frame}
going to spherical coordinates $p, \theta, \psi$ (axis along $\vr$):

\begin{equation*}
= \frac{1}{2(2\pi)^3} \int_0^\infty \dip \frac{p^2}{\omp} \int_0^\pi \dtheta \sin(\theta) e^{ipr\cos{\theta}}
\int_0^{2\pi} d\psi
\end{equation*}

integrating over $\psi$ and substituting $u = ipr\cos(\theta)$:

\begin{equation*}
\begin{split}
& = \frac{i}{2(2\pi)^2 r} \int_0^\infty \dip \frac{p}{\omp} \int_{ipr}^{-ipr} du e^{u} \\
& = \frac{i}{2(2\pi)^2 r} \int_0^\infty \dip \frac{p}{\omp} \big\{ e^{-ipr} - e^{ipr} \big\}
\end{split}
\end{equation*}

\end{frame}


\begin{frame}

juggling signs, using $\omp = \sqrtpm = \ommp$:

\begin{equation*}
\begin{split}
& = \frac{i}{2(2\pi)^2 r} \int_0^\infty \dip \frac{p}{\omp} \big\{ e^{-ipr} - e^{ipr} \big\} \\
& = \frac{-i}{2(2\pi)^2 r} \left\{ \int_0^\infty \dip \frac{-p e^{-ipr}}{\ommp} + \int_0^\infty \dip \frac{p e^{ipr}}{\omp} \right\} \\
& = \frac{-i}{2(2\pi)^2 r} \left\{ -\int_0^{-\infty} \dip \frac{p e^{ipr}}{\omp} + \int_0^\infty \dip \frac{p e^{ipr}}{\omp} \right\} \\
& = \frac{-i}{2(2\pi)^2 r} \intii \dip \frac{p e^{ipr}}{\omp}
\end{split}
\end{equation*}
\end{frame}


\begin{frame}
For the U-shaped contour we define three pieces ($\delta > 0$):
\begin{itemize}
\item $C_1$: straight line $p = i\rho - \delta$ with $\int_{C_1} \dip = i \int_\infty^m \drho$
\item $C_2$: half circle $p = im + \delta e^{i\theta}$ with $\int_{C_2} \dip = i\delta \int_\pi^{2\pi} \dtheta e^{i\theta}$
\item $C_3$: straight line $p = i\rho + \delta$ with $\int_{C_3} \dip = i \int_m^\infty \drho$
\end{itemize}
\end{frame}


\begin{frame}

With $\sqrt{m^2 - \rho^2 - 2i\rho\delta} \to -i \sqrt{\rho^2 - m^2}$:

\begin{equation*}
\begin{split}
\limdel \int_{C_1} \dip \integrand &= \limdel i \int_\infty^m \drho \frac{(i\rho - \delta)e^{i(i\rho - \delta)r}}{\sqrt{m^2 - \rho^2 - 2i\rho\delta}} \\
&= i \int_m^\infty \drho \frac{\rho e^{-\rho r}}{\sqrt{\rho^2 - m^2}}
\end{split}
\end{equation*}

With $\sqrt{m^2 - \rho^2 + 2i\rho\delta} \to i \sqrt{\rho^2 - m^2}$:

\begin{equation*}
\begin{split}
\limdel \int_{C_3} \dip \integrand &= \limdel i \int_m^\infty \drho \frac{(i\rho + \delta)e^{i(i\rho + \delta)r}}{\sqrt{m^2 - \rho^2 + 2i\rho\delta}} \\
&= i \int_m^\infty \drho \frac{\rho e^{-\rho r}}{\sqrt{\rho^2 - m^2}}
\end{split}
\end{equation*}
\end{frame}


\begin{frame}
The $\int_{C_2}$ integral vanishes for $\delta \to 0$ like $\sqrt{\delta}$:

\begin{equation*}
\left| \int_{C_2} \dip \integrand \right| \leq \delta \int_\pi^{2\pi} \dtheta
\frac{(\mathrm{bounded})}{\sqrt{2m\delta}} \propto \sqrt{\delta} \to 0
\end{equation*}

So finally, using this contour, we have:

\begin{equation*}
\begin{split}
\phixy &= \frac{-i}{2(2\pi)^2r} \integral \\
       &= \frac{1}{(2\pi)^2 r} \int_m^\infty \drho \frac{\rho e^{-\rho r}}{\sqrt{\rho^2 - m^2}}
\end{split}
\end{equation*}
\end{frame}


\begin{frame}
Approximate solution for $r \gg \frac{1}{m}$:

\begin{equation*}
\begin{split}
&\cintegral = \int_0^\infty \drho \frac{(\rho + m)e^{-(\rho + m)r}}{\sqrt{\rho(\rho + 2m)}} \\
&\approx e^{-mr} \int_0^\infty \drho \frac{m e^{-\rho r}}{\sqrt{2m\rho}}
= \sqrt{\frac{m}{2}} e^{-mr} \int_0^\infty \drho \frac{e^{-\rho r}}{\sqrt{\rho}} \\
&= \sqrt{\frac{m}{2}} e^{-mr} 2 \int_0^\infty e^{-ru^2} \du = \sqrt{\frac{\pi m}{2r}} e^{-mr}
\end{split}
\end{equation*}

Therefore:
\begin{equation*}
\phixy \approx \sqrt{\frac{m}{32\pi^3 r^3}} e^{-mr}
\end{equation*}

\end{frame}


\begin{frame}
Exact solution with the modified Bessel function $\mathrm{K}_1$:

\begin{equation*}
\begin{split}
&\cintegral =
\int_m^\infty \drho \frac{\frac{\rho}{m} e^{-\frac{\rho}{m} mr}}{\sqrt{(\frac{\rho}{m})^2 - 1}} \\
&= m \int_1^\infty \frac{u e^{-umr}}{\sqrt{u^2 - 1}} \du = m \int_1^\infty e^{-umr} \left( \di{u} \sqrt{u^2 - 1} \right) \du \\
&= -m \int_1^\infty \left( \di{u} e^{-umr} \right) \sqrt{u^2 - 1} \du = m(mr) \int_1^\infty e^{-u(mr)} \sqrt{u^2 - 1} \du \\
&= m \  \mathrm{K}_1(mr)
\end{split}
\end{equation*}

For the last equality, see \url{http://dlmf.nist.gov/10.32\#E8}.

Therefore:

\begin{equation*}
\phixy = \frac{m}{(2\pi)^2 r} \  \mathrm{K}_1(mr) \quad \textrm{for} \quad x^0 = y^0,\ mr > 0
\end{equation*}
\end{frame}


\begin{frame}
Plugging a 1-dimensionally smeared field $\phit(x) = N \int_{-\infty}^{+\infty} \dxp e^{-\frac{(x' - x)^2}{2 \sigma^2}} \phi(x')$
with $N = (\sigma \sqrt{2\pi})^{-1}$ into a bilinear form $B$ with $B(\phi(x), \phi(y)) = F(x - y)$:

\begin{equation*}
B(\phit(x), \phit(y)) = N^2 \int \dxp \dyp e^{-\frac{(x' - x)^2}{2 \sigma^2}} e^{-\frac{(y' - y)^2}{2 \sigma^2}} F(x' - y')
\end{equation*}

writing $\xi = x' - x$, $r = x - y$, $r' = x' - y'$, $\delta = r' - r$, $\sigmat = \sqrt{2} \sigma$:

\begin{equation*}
\begin{split}
&B(\phit(x), \phit(y)) = N^2 \int \dxi \drp e^{-\frac{\xi^2}{2 \sigma^2}} e^{-\frac{(\xi - (r' - r))^2}{2 \sigma^2}} F(r') \\
&= N^2 \int \dxi \drp e^{-\sigma^{-2} \left\{ (\xi - \frac{\delta}{2})^2 - \frac{\delta^2}{4} \right\}} F(r')
= N^2 \int \dxi \drp e^{-\frac{\xi^2}{\sigma^2}} e^{-\frac{\delta^2}{4 \sigma^2}} F(r') \\
&= \sigma \sqrt{\pi} N^2 \int \drp e^{-\frac{\delta^2}{2 \sigmat^2}} F(r')
= (\sigmat \sqrt{2\pi})^{-1} \int \drp e^{-\frac{(r' - r)^2}{2 \sigmat^2}} F(r')
= \frac{ \int \drp e^{-\frac{(r' - r)^2}{2 \sigmat^2}} F(r') }{ \int \drp e^{-\frac{(r' - r)^2}{2 \sigmat^2}} }
\end{split}
\end{equation*}
\end{frame}


\begin{frame}
For $B(\phi(x), \phi(y)) = F(r) = e^{ipr}$ we get:

\begin{equation*}
\begin{split}
&B(\phit(x), \phit(y)) = (\sigmat \sqrt{2\pi})^{-1} \intii \drp e^{-\frac{(r' - r)^2}{2 \sigmat^2}} e^{ipr'} \\
&= (\sigmat \sqrt{2\pi})^{-1} \intii \drp e^{-\frac{(r' - r - ip\sigmat^2)^2}{2 \sigmat^2}} e^{ipr} e^{-\frac{\sigmat^2 p^2}{2}}
\end{split}
\end{equation*}

Shifting $r' \to r'  + r + ip\sigmat^2$ for $|r'| < R$ and $r' \to r' + r$ for $|r'| \gg R$ and letting $R \to \infty$, we get:

\begin{equation*}
\begin{split}
&B(\phit(x), \phit(y)) = (\sigmat \sqrt{2\pi})^{-1} \intii \drp e^{-\frac{{r'}^2}{2 \sigmat^2}} e^{ipr} e^{-\frac{\sigmat^2 p^2}{2}} \\
&= e^{ipr} e^{-\frac{\sigmat^2 p^2}{2}} = e^{ipr} e^{-(\sigma p)^2}
\end{split}
\end{equation*}
\end{frame}


\end{document}
